% =====================================================================
%  บทคัดย่อภาษาไทย
%  ข้อมูลหัวข้อ ชื่อผู้เขียน อาจารย์ที่ปรึกษา ฯลฯ ดึงมาจาก main.tex อัตโนมัติ
%  คำสำคัญ กำหนดด้วย \keywordsth{...} ใน main.tex
% =====================================================================
\begin{abstractth}
งานวิจัยนี้มีวัตถุประสงค์เพื่อพัฒนาระบบผู้ช่วยสอนที่ใช้ปัญญาประดิษฐ์เชิงสร้างสรรค์ (Generative AI) สำหรับรายวิชาการเขียนโปรแกรมเบื้องต้น และเพื่อประเมินประสิทธิภาพของระบบในการช่วยให้ผู้เรียนเข้าใจแนวคิดการเขียนโปรแกรมและแก้ไขข้อผิดพลาดของโค้ดได้ด้วยตนเอง ระบบที่พัฒนาขึ้นใช้แบบจำลองภาษาขนาดใหญ่ (Large Language Model: LLM) ร่วมกับเทคนิคการค้นคืนข้อมูลเสริมการสร้าง (Retrieval-Augmented Generation: RAG) เพื่ออ้างอิงเนื้อหาจากเอกสารประกอบการสอนของรายวิชา

กลุ่มตัวอย่างคือนักศึกษาชั้นปีที่ 1 จำนวน 60 คน แบ่งเป็นกลุ่มทดลองและกลุ่มควบคุมกลุ่มละ 30 คน ผลการวิจัยพบว่า กลุ่มทดลองมีคะแนนผลสัมฤทธิ์ทางการเรียนหลังเรียนสูงกว่ากลุ่มควบคุมอย่างมีนัยสำคัญทางสถิติที่ระดับ .05 และมีความพึงพอใจต่อระบบในระดับมากที่สุด (ค่าเฉลี่ย 4.52 จากคะแนนเต็ม 5)

\textit{(ข้อความนี้เป็นตัวอย่าง บทคัดย่อควรมีความยาวไม่เกิน 1 หน้า ประกอบด้วยวัตถุประสงค์ วิธีการ ผลการวิจัย และข้อสรุปโดยย่อ)}
\end{abstractth}
